\addcontentsline{toc}{subsubsection}{Solution 32. The Skew-Symmetric Parallelogram Squeeze}

\begin{dlsummary}[title={\bfseries Solution 32. The Skew-Symmetric Parallelogram Squeeze},breakable]

\begin{proof}
\textbf{Step 1. Obtain a quadratic lower bound.}

Comparing the outer expressions in the given chain gives
\[
2xy+x^2f(y)-y^2f(x)\leq f(x)+y^2.
\]
Setting $y=x$ cancels the skew-symmetric terms, so
\[
2x^2\leq f(x)+x^2,
\qquad\text{hence}\qquad
f(x)\geq x^2\quad(x\in\mathbb R). \tag{1}
\]
Define $g(x)=f(x)-x^2$. Then $g(x)\geq0$ for all real $x$.

\medskip
\noindent
\textbf{Step 2. Establish midpoint concavity of the residual.}

The right-hand inequality and the parallelogram identity give
\[
\frac{(x+y)^2+(x-y)^2}{2}
+\frac{g(x+y)+g(x-y)}2
\leq x^2+y^2+g(x).
\]
After cancellation,
\[
g(x+y)+g(x-y)\leq2g(x)
\qquad(x,y\in\mathbb R). \tag{2}
\]

\medskip
\noindent
\textbf{Step 3. Prove that $g$ is constant.}

Fix $a,b\in\mathbb R$, and put
\[
x_k=a+k(b-a)\quad(k\geq0),
\qquad
\Delta_k=g(x_{k+1})-g(x_k).
\]
For each $k\geq1$, applying (2) with $x=x_k$ and $y=b-a$ yields
\[
g(x_{k+1})+g(x_{k-1})\leq2g(x_k),
\]
so $\Delta_k\leq\Delta_{k-1}$. Thus $\Delta_k\leq\Delta_0=g(b)-g(a)$
for all $k\geq0$. For every positive integer $n$, nonnegativity of $g$
implies
\[
-g(a)\leq g(x_n)-g(a)
=\sum_{k=0}^{n-1}\Delta_k
\leq n\bigl(g(b)-g(a)\bigr).
\]
Letting $n\to\infty$ gives $g(b)\geq g(a)$. Interchanging $a$ and $b$
shows $g(a)=g(b)$, so $g(x)=d$ for some constant $d\geq0$. Therefore
\[
f(x)=x^2+d. \tag{3}
\]

\medskip
\noindent
\textbf{Step 4. Eliminate the constant shift.}

Substituting (3) into the outer inequality gives
\[
(x-y)^2+d(1-x^2+y^2)\geq0
\qquad(x,y\in\mathbb R).
\]
Taking $y=x-1$ reduces this to
\[
1+d(2-2x)\geq0\qquad(x\in\mathbb R).
\]
If $d>0$, this fails for sufficiently large $x$. Hence $d=0$ and
$f(x)=x^2$ for all real $x$.

For $f(x)=x^2$, the middle and right-hand expressions are both
$x^2+y^2$, while the left-hand expression is $2xy\leq x^2+y^2$.
Thus this function satisfies the stated inequalities and is the unique
solution.
\end{proof}

\end{dlsummary}